% L15a strategy deck: the gamma-by-sector picking rule.
% One short deck per rule on the strategy menu, read next to the rule's code file
% (strategies/picking/gamma_by_sector.jl). Written by make_strategy_decks.py (session scratchpad, 2026-11-27) and then
% kept by hand: edit this file, not the generator. Build with make in this folder: LaTeX runs from strategies/,
% so paths are relative to it.
\documentclass[aspectratio=169,11pt]{beamer}
\usepackage{vnslides}

\newcommand{\repo}{https://github.com/varnerlab/CHEME-5660-CourseRepository-Fall-2026/blob/main/lectures/week-15/L15a}

\title{\texorpdfstring{{\vnheavy L15a}}{L15a}: Gamma by Sector}
\subtitle{Picking Rule, CHEME 5660 Strategy Menu, Cornell University}
\author{Copyright Jeffrey Varner 2026. All Rights Reserved}

\begin{document}

\vntitlepage

\begin{frame}{Disclaimer and Risks}
  \vspace{0.6em}
  \textbf{\ink{Training only.}} This content is for training and informational
  purposes only. The teaching team does not make, give, or endorse any offer or
  solicitation to buy or sell securities or derivative products, or provide
  investment or trading advice or strategy.

  \vspace{1.1em}
  \textbf{\ink{Risk.}} Trading involves risk and is generally not recommended for
  those with limited resources, experience, or a low risk tolerance. Past
  performance does not guarantee future success.

  \vspace{1.1em}
  \textbf{\ink{Responsibility.}} You are responsible for your investment decisions
  based on your financial circumstances, objectives, risk tolerance, and
  liquidity needs.
\end{frame}

\begin{frame}{The Idea}
  \vspace{0.3em}
  The gamma-by-sector rule picks the basket from the preference weights, three tickers in each sector:
  \begin{itemize}\setlength{\itemsep}{0.5em}
    \item \hilite{Universe}: 491 stocks in 11 sectors and 14 exchange-traded funds (ETFs), without \texttt{SPY} and the Treasury bill fund \texttt{SGOV}. The ETFs count as one more sector.
    \item \hilite{Basket}: the three tickers with the highest preference weights in each of the 12 sectors, 36 tickers in all.
    \item \hilite{Reason}: taking the same number from each sector spreads the basket across the economy.
  \end{itemize}

  \vspace{0.5em}
  The engine runs the rule at each decision: on the start day and whenever a trigger fires. Between decisions, the basket stays the same.
\end{frame}

\begin{frame}{The Preference Weight}
  \vspace{0.3em}
  At each close $t$, the engine scores every ticker $i$ with its preference weight $\gamma_{i}(t)$:
  {\setlength{\jot}{3pt}
  \begin{align*}
    \tilde{g}_{i,t} &= \alpha_{i}(t)+\beta_{i}(t)\,\tilde{g}_{M,t} && \text{the expected growth rate (1/yr)}\\
    \gamma_{i}(t) &= \tanh\left(h\,\frac{\tilde{g}_{i,t}}{\max\left(|\beta_{i}(t)|,\,\beta_{\text{floor}}\right)^{\xi_{t}}}\right)
      && h=1\text{ year},\ \beta_{\text{floor}}=0.1
  \end{align*}}%
  where the single index model gives $\alpha_{i}(t)$ and $\beta_{i}(t)$, $\tilde{g}_{M,t}$ is the 252-day
  exponential moving average of market growth, and the market state $\xi_{t}$ is positive when the market falls.

  \vspace{0.5em}
  The preference weight lies between $-1$ and $1$, with the sign of $\tilde{g}_{i,t}$.

  \vspace{0.5em}
  With fixed estimates, $\alpha_{i}(t)$ and $\beta_{i}(t)$ keep their values from the 2016 to 2024 fit. With
  updated estimates, the engine updates them at every close. The \href{\repo/strategies/slides/CHEME-5660-L15a-Strategy-Sector-Factor-Fall-2026.pdf}{sector factor estimates}
  keep the fixed values and add the sector growth that the market does not explain, $\hat{\delta}_{i}\tilde{f}_{s,-i}(t)$, to $\tilde{g}_{i,t}$.
\end{frame}

\begin{frame}{How the Market State Tilts the Preference Weight}
  \vspace{0.3em}
  The market state $\xi_{t}$ sets how the market sensitivity scales the expected growth rate:
  \begin{center}
    \includegraphics{slides/figs/Fig-L15a-Deck-Tilt.pdf}
  \end{center}
  At the same positive expected growth rate, the preference weight favors a low market sensitivity when the
  market falls and a high one when the market rises.
\end{frame}

\begin{frame}{The Rule}
  \vspace{0.3em}
  On decision day $t$, split the universe into its sectors $\mathcal{S}_{1},\dots,\mathcal{S}_{12}$:
  \begin{enumerate}\setlength{\itemsep}{0.5em}
    \item In each sector $\mathcal{S}_{j}$, rank the tickers by $\gamma_{i}(t)$, from highest to lowest.
    \item Take the first three tickers of each ranking. A sector with three tickers or fewer gives all of them.
    \item The basket $\mathcal{B}(t)$ holds the tickers taken in every sector.
  \end{enumerate}

  \vspace{0.5em}
  The order of the tickers in the universe breaks ties. The rule takes the top three tickers of a sector
  even when their preference weights are negative. The weighting rule then decides how much of each ticker to hold.
\end{frame}

\begin{frame}{The Basket of a Made-Up Market}
  \vspace{0.3em}
  In a made-up market of 12 sectors, the rule takes the top three of each sector:
  \begin{center}
    \includegraphics{slides/figs/Fig-L15a-Deck-Picks-Sector.pdf}
  \end{center}
  In Health and Real estate, the third ticker taken has a negative preference weight. The random-by-sector rule, on the
  right, draws three tickers of each sector at random.
\end{frame}

\begin{frame}{Reading the Code}
  \vspace{0.3em}
  The rule's code file is
  \href{\repo/strategies/picking/gamma_by_sector.jl}{\texttt{strategies/picking/gamma\_by\_sector.jl}}:
  \begin{itemize}\setlength{\itemsep}{0.5em}
    \item \texttt{pick\_gamma\_by\_sector(tickers, sectors, gamma; count)} takes the \texttt{count} tickers with the highest preference weights in each sector.
    \item \texttt{MyL15aGammaSectorScreen} is the type the engine calls. The rule has no settings of its own.
    \item \texttt{MPC.pick(screen, md, t; $\gamma$\_t, N, ...)} runs at each decision with the day's preference weights \texttt{$\gamma$\_t} and \texttt{N~=~3} tickers per sector, and returns the basket's columns in the market data.
  \end{itemize}

  \vspace{0.5em}
  The \href{\repo/strategies/CHEME-5660-L15a-Strategy-Menu-Fall-2026.ipynb}{strategy menu} lists every rule, and the \href{\repo/engine/CHEME-5660-L15a-Trading-Engine-Fall-2026.ipynb}{trading engine} describes the preference weight and the triggers.
\end{frame}

\end{document}
